\begin{table}[htbp]
\centering
\small
\caption{Asset-Disaggregated Capital Formation and Capacity Utilization (1968--1975): Empirical Falsification of Austrian Malinvestment (H2)}
\label{tab:austrian_falsification_capital}
\resizebox{\textwidth}{!}{%
\begin{tabular}{lcccccccc}
\toprule
& \multicolumn{2}{c}{\textbf{Pre-UP Administration}} & \textbf{Baseline} & \multicolumn{3}{c}{\textbf{Unidad Popular Administration}} & \multicolumn{2}{c}{\textbf{Post-Coup Shock Adjustment}} \\
\cmidrule(lr){2-3} \cmidrule(lr){4-4} \cmidrule(lr){5-7} \cmidrule(lr){8-9}
\textbf{Variable / Asset Class} & \textbf{1968} & \textbf{1969} & \textbf{1970} & \textbf{1971} & \textbf{1972} & \textbf{1973} & \textbf{1974} & \textbf{1975} \\
\midrule
\multicolumn{9}{l}{\textbf{A. Real Gross Fixed Investment by Asset Class ($I_t$, Millions Real 2003 CLP)}} \\[0.2em]
Machinery \& Equipment ($I_{\text{ME}}$) & 3.14 & 3.09 & 3.26 & 2.95 & 2.27 & 2.53 & 2.55 & 2.60 \\
Non-Residential Structures ($I_{\text{NRC}}$) & 3.05 & 3.21 & 3.68 & 3.34 & 2.64 & 2.37 & 3.20 & 2.02 \\
Residential Housing ($I_{\text{RC}}$) & 2.43 & 2.75 & 2.69 & 3.12 & 2.61 & 2.17 & 2.66 & 1.99 \\
Total Gross Investment ($I_{\text{Total}}$) & 8.62 & 9.05 & 9.64 & 9.41 & 7.52 & 7.07 & 8.42 & 6.61 \\
\midrule
\multicolumn{9}{l}{\textbf{B. Cumulative Investment Growth Relative to 1970 Baseline (\% $\Delta$)}} \\[0.2em]
\% $\Delta I_{\text{ME}}$ (Machinery) & -3.7\% & -5.3\% & 0.0\% & \textbf{-9.5\%} & \textbf{-30.5\%} & \textbf{-22.4\%} & -21.6\% & -20.4\% \\
\% $\Delta I_{\text{NRC}}$ (Productive Structures) & -17.2\% & -12.8\% & 0.0\% & \textbf{-9.2\%} & \textbf{-28.3\%} & \textbf{-35.6\%} & -13.0\% & -45.2\% \\
\% $\Delta I_{\text{RC}}$ (Residential Housing) & -9.9\% & +2.2\% & 0.0\% & \textbf{+15.7\%} & -3.0\% & -19.5\% & -1.2\% & -25.9\% \\
\midrule
\multicolumn{9}{l}{\textbf{C. Structural Capacity Utilization and Relative Price Dynamics}} \\[0.2em]
Capacity Utilization Index ($\mu_t$) & 0.9366 & 0.9410 & 0.9278 & \textbf{0.9844} & \textbf{0.9599} & 0.8946 & 0.8882 & \textbf{0.7662} \\
Capital Deflator ($P^K$, 2003 = 1.0) & 0.000343 & 0.000461 & 0.000653 & 0.000839 & 0.001746 & 0.008916 & 0.079644 & 0.415686 \\
Gross Capital Stock ($K^g_{\text{Total}}$, M 2003 CLP) & 154.5 & 160.3 & 166.6 & 172.6 & 176.5 & 179.9 & 184.5 & 187.1 \\
Net Capital Stock ($K^n_{\text{Total}}$, M 2003 CLP) & 95.8 & 99.4 & 103.3 & 107.0 & 109.5 & 111.6 & 114.4 & 116.1 \\
\bottomrule
\multicolumn{9}{p{17.2cm}}{\footnotesize \textbf{Sources and Notes:} Capital formation series ($I_t$), gross stocks ($K^g$), net stocks ($K^n$), and capital goods deflator ($P^K$) sourced from the harmonized perpetual inventory dataset of Chapter~2. Structural capacity utilization ($\mu_t$) estimated via the Cointegrating Multivariate Polynomial Regression (CMPR IM-OLS) framework of Chapter~2. Austrian Business Cycle Theory ($H_2$) predicts that credit expansion artificially lowers interest rates, inducing an over-investment boom in higher-order capital goods (machinery $I_{\text{ME}}$ and non-residential productive structures $I_{\text{NRC}}$). The audited data directly refutes this hypothesis: real productive investment collapsed in 1971 ($-9.5\%$) and 1972 ($-30.5\%$), while capacity utilization reached a historical peak of $98.44\%$, demonstrating an import-constrained investment squeeze rather than Austrian malinvestment.}
\end{tabular}%
}
\end{table}
